\lesson{11}{Mar 05 2022 Sat (01:03:47)}{Discriminant of a quadratic}{Unit 4}

\begin{review}[Quadratic Equation]
  \label{rev:quadratic_equation}

  \subsubsection{Introduction}
  \label{sub_sub_sec:introduction}

  A \textbf{Quadratic Equation} is when the variable is squared. So we can have
  something called a \textbf{Quadratic Expression}, which looks something like
  this: $x^{2} + 8x + 15$. Most of the time, quadratic equations are written in
  the form of: $x^{2} \pm ax \pm b = 0$. But, you may see them sometimes written
  in this form: $ax^{2} + bx + c = 0$. That form is called the
  \textbf{Standard Form}.

  We can do a lot of things with \textbf{Quadratic Expressions}. For one thing,
  we can:

  \begin{enumerate}
    \label{enum:quadratic_equation}
  
    \item Factor them (most of the time). The equation above will be factored
      to $(x + 3)(x + 5)$.
    \item Graph them. We take the expression and re-write it as a function.
      $f(x) = x^{2} + 8x + 15$.

      \begin{figure}[H]
        \centering
        \incfig{quadratic-equation}
        \label{fig:quadratic_equation}
      \end{figure}
  \end{enumerate}

  % subsubsection introduction (end)

  \subsubsection{Solving a Quadratic Equation}
  \label{sub_sub_sec:solving_a_quadratic_equation}

  Let's solve for $x$ : $2x^{2} + 3x - 2 = 0$.

  For this, we need the quadratic equation:
  $x = \frac{-b \pm \sqrt{b^{2} - 4ac}}{2a}$

  \begin{equation}
    \label{eqt:solving_a_quadratic_equation}
    \begin{split}
      x &= \frac{-b \pm \sqrt{b^{2} - 4ac}}{2a} \\
      2x^{2} + 3x - 2 = 0 \implies x &= \frac{-3 \pm \sqrt{3^{2} - 4(2)(-2)}}{2(2)} \\
      &= \frac{-3 \pm \sqrt{9 + 16}}{4} \\
      &= \frac{-3 \pm \sqrt{25}}{4} \\
      &= \frac{-3 += 5}{4} \\
      x = \frac{1}{2} \qquad \qquad &\qquad \qquad x = -2 \\
    \end{split}
  .\end{equation}

  % subsubsection solving_a_quadratic_equation (end)
\end{review}

\begin{definition}[Discriminant]
  \label{def:discriminant}

  The expression $b^{2} - 4ac$ that appears under the square root sign in the
  quadratic equation (\ref{rev:quadratic_equation}) is called the
  \textbf{Discriminant} of the quadratic equation.
\end{definition}

The solutions of the quadratic equation $ax^{2} + bx + c = 0$ are given by the
quadratic formula:

The expression under the square root, $b^{2} - 4ac$, is the
\textbf{Discriminant} (\ref{def:discriminant}) of the quadratic equation.

The value of the discriminant determines the number of real solutions for
$ax^{2} + bx + c = 0$.

\begin{enumerate}
  \label{enum:number_of_real_solutions}

  \item if $b^{2} - 4ac > 0$, there are $2$ real solutions to
    $ax^{2} + bx + c = 0$
  \item if $b^{2} - 4ac = 0$, there is $1$ real solution to
    $ax^{2} + bx + c = 0$
  \item if $b^{2} - 4ac < 0$, there are no real solution to
    $ax^{2} + bx + c = 0$
\end{enumerate}

\begin{example}
  \label{exm:find_discriminant}

  For the quadratic equation $5x^{2} - 9x + 1 = 0$, we have
  $a = 5 \qquad b = -9 \qquad c = 1$. So, the value of the discriminant is as
  the follows:

  \begin{equation}
    \label{eqt:find_discriminant}
    \begin{split}
      b^{2} - 4ac = (-9)^{2} - 4(5)(1) = 81 - 21 = 61
    \end{split}
  .\end{equation}

  Since the discriminant is $> 0$, $5x^{2} - 9x + 1 = 0$ has $2$ real solutions.
\end{example}

\newpage
